% !TEX root = ../notes.tex

\documentclass[../notes.tex]{subfiles}

\begin{document}

\section{March 11}
Here we go.

\subsection{Characteristic Zero to Positive Characteristic, Again}
For today, for primes $p$, we set $k_p\coloneqq\overline\FF_p$ and $k_0\coloneqq\CC$. We then choose some characteristic $p$, and we let $G_p$ be a connected reductive group over $k_p$, we let $U_p\subseteq G_p(k_p)$ is the set of unipotent elements, and $\op{Conj}(U_p)$ is the set of unipotent conjugacy classes. Last time, we showed that there is an embedding $\op{Conj}(U_0)\to\op{Conj}(U_p)$; note that these are finite sets.
\begin{remark}
	One has that $\op{Conj}(U_p)$ can be stratified into a disjoint union of locally closed subsets (of $U_p$) given by the $H^{\overline\omega}$s, and these locally closed subsets have the same indexing set as for $\op{Conj}(U_0)$.
\end{remark}
Now, for any of our $\mc P\subseteq U_p$ and finite field $\FF_q\subseteq k_p$, one has that $\left|\mc P(\FF_q)\right|$ is independent of $p$ (as a function of $q$).
\begin{example}
	Consider $\op{Sp}_4(\FF_q)$. This has four pieces for odd $q$.
	\begin{itemize}
		\item There is the identity (with $1$ element).
		\item There is a class with $q^4-1$ elements.
		\item The subregular unipotent elements split into two classes, which have $\frac12q(q-1)\left(q^4-1\right)$ elements and $\frac12q(q+1)\left(q^4-1\right)$ elements. (This rejoins as a single conjugacy class over $\overline\FF_q$.)
		\item The regular unipotent elements split into two classes, each with $\frac12q^2\left(q^2-1\right)\left(q^4+1\right)$ elements.
	\end{itemize}
	For $p=2$, the subregular elements split into two conjugacy classes (even over $\overline\FF_2$) which have a different number of elements than above (namely, $\left(q^2-1\right)\left(q^4-1\right)$ and $q^4-1$), but the union of all the subregular unipotent elements is the same!
\end{example}
\begin{remark}
	Steinberg has shown that the total number of unipotent elements is a power of $q$.
\end{remark}
The point of this discussion is that we may define the following.
\begin{notation}
	We define $\op{Conj}(U_*)$ to be the quotient of $\bigsqcup_{p>0}\op{Conj}(U_p)$ by the equivalence relation $\sim$ given by $x_0\sim x_p$ for any $x_0\in\op{Conj}(U_0)$.
\end{notation}
\begin{example}
	For classical groups, it turns out that $\op{Conj}(U_*)=\op{Conj}(U_2)$.
\end{example}
\begin{example}
	For $E_8$, one has that $\op{Conj}(U_*)=(\op{Conj}(U_2)\sqcup\op{Conj}(U_3))/\op{Conj}(U_0)$. Here, $\#\op{Conj}(U_2)=74$ and $\#\op{Conj}(U_3)=71$, so $\#\op{Conj}(U_*)=74$.
\end{example}
We are going to use $\op{Conj}(U_*)$ to index the strata of $G$ discussed last time.

Let's use the Springer representation to give another embedding $\op{Conj}(U_0)\to\op{Conj}(U_p)$. Let's start with symplectic types.
\begin{notation}
	Recall that the irreducible representations of the Weyl group in type $C_n$ is given by a sequence $(a_i)$ with $a_0\le a_2\le\cdots\le a_{2m}$ and $a_1\le a_3\le\cdots\le a_{2m-1}$ and total sum $\sum_ia_i=n$; two such sequences are the same if we prepend by two $0$s. These representations are found by inducing from the two copies of $\op{GL}_m$ in a parabolic. We might label this set by $\op{Irr}C_n$.
\end{notation}
\begin{remark}
	Professor Lusztig labeled this set as $\op{Irr}B_n$ because the Weyl group in types $B_n$ and $C_n$ are the same. We will continue to use $C_n$ to avoid confusion.
\end{remark}
Here are some alternative parameterizations.
\begin{notation}
	We define $\op{Sym}C_n$ to be given by sequences $(b_i)$, where we replace the $\le$s requirements with $<$ and replace the sum condition to summing to $n+m^2$, with the same equivalence. Note that there is a bijection $\op{Irr}C_n\to\op{Sym}C_n$ by taking $(a_i)$ to $(a_i+\floor{i/2})_i$. The equivalence classes of increasing sequences $(b_i)$ parameterize the ``special'' representations of the Weyl group.
\end{notation}
The notion of ``special'' representation was defined in a paper of Lusztig in 1978.
\begin{notation}
	We define $\op{Sym}'C_n$ to be given by sequences $(c_i)$, where we replace the $\le$s requirements with $\le2+$ and replace the sum condition to summing to $n+2m^2+m$. One should also adjust the equivalence to be $(c_i)\sim(0,1,c_0+2,c_1+2,\ldots)$. Note that there is a bijection $\op{Sym}C_n\to\op{Sym}'C_n$ by taking $(b_i)$ to $(b_i+\floor{(i+1)/2})_i$. The equivalence classes of increasing sequences $(c_i)$ parameterizes the representations of $W$ which correspond to unipotent classes under the Springer correspondence for $p\ne2$.
\end{notation}
\begin{notation}
	We define $\op{Sym}''C_n$ to be given by sequences $(d_i)$, where we replace the $\le$s requirements with $\le4+$ and replace the sum condition to summing to $n+4m^2+m$; we also need $d_1\ge2$. One should also adjust the equivalence to be $(d_i)\sim(0,2,c_0+4,c_1+4,\ldots)$. Note that there is a bijection $\op{Sym}'C_n\to\op{Sym}''C_n$ by taking $(c_i)$ to $(c_i+i)_i$. The equivalence classes of increasing sequences $(d_i)$ parameterizes the representations of $W$ which correspond to unipotent classes under the Springer correspondence for $p=2$.
\end{notation}
The statements about the unipotent classes were worked out by Lusztig (and Spaltenstein for $p=2$) in the 1980s. We remark that these parameterizations automatically provide embeddings $\op{Conj}(U_p)\to\op{Conj}(U_2)$ for $p\ne2$.
\begin{remark}
	One can do something similar in other types. In particular, basically everything is the same in the odd orthogonal types $B_n$, except $\op{Sym}'$ changes slightly. In type $D_n$, all the definitions are the same, except one removes $a_0$s from $\op{Irr}D_n$, removes $b_0s$ from $\op{Sym}D_n$ (and sums to $n+m^2-m$), removes $c_0$s from $\op{Sym}'D_n$ (and sums to $n+2m^2-2m$), and removes $d_0$s from $\op{Sym}''D_n$ (and sums to $n+4m^2-4m$). Importantly, one also needs to adjust our equivalence so that $(a_1,a_2,a_3,a_4,\ldots)\sim(a_2,a_1,a_4,a_3,\ldots)$, and when this movement does nothing, our sets should have two copies of this sequence. The statements about the parameterizations are all the same.
\end{remark}
Here is some more notation.
\begin{notation}
	Let $\op{Irr}_*W$ be the collection of irreducible representations of $W$ appearing in the Springer correspondence. We also let $\op{Irr}_pW$ be the image of $\op{Conj}(U_p)$ in $\op{Irr}_*W$ under the Springer correspondence.
\end{notation}
\begin{remark}
	It turns out that
	\[\op{Irr}_*W=\bigcup_{\text{prime }p}\op{Irr}_pW.\]
\end{remark}
\begin{remark}
	The Springer correspondence provides a bijection $\op{Conj}(U_*)\to\op{Irr}_*W$ which restricts to isomorphisms $\op{Conj}(U_p)\to\op{Irr}_pW$ for each prime $p$.
\end{remark}
\begin{example}
	For $E_8$, one finds that $\op{Irr}W$ has $112$ elements, but $\op{Irr}_*W$ has $75$ elements.
\end{example}
The moral is that the indexing of our strata will be independent of the characteristic.

\subsection{\texorpdfstring{$2$}{2}-Special Representations}
Let $V$ and $V^*$ be dual rational vector spaces, and let $\langle-,-\rangle$ be the induced pairing. We also let $\Phi\subseteq V$ and $\Phi^\lor\subseteq V^*$ be dual root systems, and we let $W$ be the corresponding Weyl group generated by the (simple) reflections.
\begin{notation}
	For $e\in V$, we define
	\[\Phi_e\coloneqq\{\alpha\in\Phi:\langle e,\alpha^\lor\rangle\in\ZZ\},\]
	and
	\[\Phi^\lor_e\coloneqq\{\alpha^\lor\in\Phi:\alpha\in\Phi e\}.\]
	Dually, for $e\in V^*$, we may also define $\Phi_e$ and $\Phi^\lor_e$ as above (replacing $\Phi$ with $\Phi^\lor$ and adjusting equations as needed).
\end{notation}
\begin{remark}
	Now, for a pair $(e,e')$, it turns out that $(V,V^*,\Phi_e\cap\Phi_{e'},\Phi^\lor_e\cap\Phi^\lor_{e'})$ is a root system. We let its Weyl group be $W_{ee'}$.
\end{remark}
\begin{example}
	Here are the possible subgroups $W_{ee'}$ of the same rank.
	\begin{itemize}
		\item In type $A_n$, the Weyl subgroups $W_{ee'}$ are $A_n$.
		\item In types $B_n$ and $C_n$, the Weyl subgroups $W_{ee'}$ take the form $B_?\times B_?\times D_?\times D_?$.
		\item For type $D_n$, the Weyl subgroups look like $D_?\times D_?\times D_?\times D_?$.
		\item There are also lists for the exceptional types.
	\end{itemize}
\end{example}
Let's explain why this relates to $\op{Irr}_*W$.
\begin{notation}
	For an irreducible representation $E$ of $W_{ee'}$, we let $n_E$ be the smallest nonnegative integer for which $E\subseteq V^{\otimes n_E}$, where $V$ is the standard (reflection) representation of $W_{ee'}$. We have a similar definition for irreducible definitions of $W$.
\end{notation}
\begin{notation}
	Let $E$ be the special representation of $W_{ee'}$. Then its induction to $W$ has a unique irreducible component $\widetilde E$ for which $n_{\widetilde E}=n_E$, which we denote by $j^W_{W_{ee'}}(E)$.
\end{notation}
\begin{definition}[special]
	An irreducible representation $\widetilde E$ is \textit{$2$-special} if and only if it takes the form $j^W_{W_{ee'}}(E)$ for some $(e,e')\in V\times V^*$ and some special representation $E$ of $W_{ee'}$.
\end{definition}
It turns out that the $2$-special representations of $W$ are in bijection with $\op{Irr}_*W$.
\begin{remark}
	It turns out that there is another definition of a $1$-special representation, where one deals with $W_e$ instead of $W_{ee'}$. If one could define $3$-special representations, Professor Lusztig expects that they would parameterize the strata.
\end{remark}

\end{document}